Define sublinear functions, relate them to seminorms

Recall subadditivity: \( f(x + y) \leq f(x) + f(y) \) for all \( x, y \in X \)

A real function \( f : X \to \mathbb{R} \) on a vector space is sublinear if it is both
positively homogeneous and subadditive, or equivalently, if it is both positively
homogeneous and convex.

A seminorm is a subadditive function \( p : X \to \mathbb{R} \) satisfying
\( p(\alpha x) = |\alpha|x \)
for all \( \alpha \in \mathbb{R} \) and \( x \in X \).
A seminorm \( p \) is a norm if and only if \( p(x) = 0 \iff x = 0 \)

Every seminorm is sublinear and a sublinear function satisfying \( f(x) = f(-x) \) is also a seminorm.

If \( p : X \to \mathbb{R} \) is sublinear, then:
1. \( p(0) = 0 \).
2. For all \( x \) we have \( - p(x) \leq p(-x) \). Consequently \( p \) is linear if and only if
\( p(-x) = - p(x) \) for all \( x \in X \).
3. The function \( q(x) := \text{max}(p(x), p(-x)) \) is a seminorm.
4. If \( p \) is a seminorm, then \( p(x) \geq 0 \) for all \( x \).
5. If \( p \) is a seminorm, then the set \( \{ x : p(x) = 0\} \) is a linear subspace.

A nonnegative sublinear function on a tvs is:
1. Lower semicontinuous if and only if it is the gauge/Minkowski functional of an absorbing closed
convex set.
2. Continuous if and only if it is the gauge of a convex neighborhood of zero.

(Verification of linear cont. functions)
A linear functional on a tvs is continuous if and only if it is dominated by a continuous seminorm.



What are Gauges/Minkowski functionals? State some useful theorems about them

The gauge or the Minkowski functional, \( p_A \) of a subset \( A \subseteq X \) 
is defined as
\( p_A(x) = \text{inf} \{ \alpha > 0 : x \in \alpha A\} \)
where, by convention, \( \text{inf} \emptyset = \infty \) 
In other words, \( p_A(X) \) is the smallest factor by which the set \( A \) 
must be enlarged to contain the point \( x \).

The gauge satisfies the following properties:
**Lemma**  
For nonempty subsets \( B \) and \( C \) of a vector space \( X \):
1. \( p_C(x) = p_C(-x) \) for all \( x \in X \).
2. If \( C \) is symmetric, then \( p_C(x) = p_C(-x) \) for all \( x \in X \).
3. \( B \subset C \) implies \( p_C \leq p_B \).
4. If \( C \) includes a subspace \( M \), then \( p_C(x) = 0 \) for all \( x \in M \).
5. If \( C \) is star-shaped about zero, then  
   \[
   \{ x \in X : p_C(x) < 1 \} \subset C \subset \{ x \in X : p_C(x) \leq 1 \}.
   \]
6. If \( X \) is a tvs and \( C \) is closed and star-shaped about zero, then  
   \[
   C = \{ x \in X : p_C(x) \leq 1 \}.
   \]
7. If \( B \) and \( C \) are star-shaped about zero, then \( p_{B \vee C} = p_B \vee p_C \), where as usual  
   \[
   (p_B \vee p_C)(x) = \max \{ p_B(x), p_C(x) \}.
   \]


For a nonnegative function \( p : X \to \mathbb{R} \) on a vector space, we have the following:
1. \( p \) is positively homogeneous if and only if it is the gauge of an absorbing set—in which case for every subset \( A \) of \( X \) satisfying  
   \[
   \{ x \in X : p(x) < 1 \} \subset A \subset \{ x \in X : p(x) \leq 1 \}
   \]  
   we have \( p_A = p \).
2. \( p \) is sublinear if and only if it is the gauge of a convex absorbing set \( C \), in which case we may take \( C = \{ x \in X : p(x) \leq 1 \} \).

3. \( p \) is a seminorm if and only if it is the gauge of a circled convex absorbing set \( C \), in which case we may take \( C = \{ x \in X : p(x) \leq 1 \} \).
4. When \( X \) is a tvs, \( p \) is a continuous seminorm if and only if it is the gauge of a unique closed, circled and convex neighborhood \( V \) of zero, namely  
   \[
   V = \{ x \in X : p(x) \leq 1 \}.
   \]
5. When \( X \) is finite dimensional, \( p \) is a norm if and only if it is the gauge of a unique circled, convex and compact neighborhood \( V \) of zero, namely  
   \[
   V = \{ x \in X : p(x) \leq 1 \}.
   \]



What is a dual pairing and the corresponding topology?

A dual pair (or a dual system) is a pair \((X, X')\) of vector spaces together with a bilinear functional \((x, x') \mapsto \langle x, x' \rangle\), from \(X \times X'\) to \(\mathbb{R}\), that separates the points of \(X\) and \(X'\). That is:
1. The mapping \(x' \mapsto \langle x, x' \rangle\) is linear for each \(x \in X\).
2. The mapping \(x \mapsto \langle x, x' \rangle\) is linear for each \(x' \in X'\).
3. If \(\langle x, x' \rangle = 0\) for each \(x' \in X'\), then \(x = 0\).
4. If \(\langle x, x' \rangle = 0\) for each \(x \in X\), then \(x' = 0\).

3-4 mean that they separate points of each other (total).

Since we can consider X to be a vector subspace of \( \mathbb{R}^{X'} \), \( X \)  inherits the product topology of \( \mathbb{R}^X \). 
This topology is referred to as the weak topology on \( X \) and is denoted \( \sigma(X, X') \), 
or simply \( w \). Since the product topology on \( \mathbb{R}^{X'} \) is a locally
convex Hausdorff topology, the weak topology \( \sigma(X, X') \) is likewise Hausdorff and locally convex.
As \( x_{\alpha} \to^{w} x \iff \langle x_{\alpha}, x' \rangle \to \langle x, x' \rangle \) in \( \mathbb{R} \)
this is also called the topology of pointwise convergence on \( X' \).

A family of seminorms that generates the weak
topology \( \sigma(X, X') \) is \( \{p : x' : x' \in X'\} \), where
\( p_{x'}(x) = \|\langle x, x' \rangle\| \)

The analogous construction for \( \sigma(X', X) \) is called the \( w^\ast \)-topology
The central strength with this construction is that all possible topologies are identified by certain choices of the dual
(which is not unique in the infinite dimensional case)

If \( \langle X, X' \rangle \) is a dual pair, then
the topological dual of the TVS \( (X, \sigma(X, X')) \) is \( X' \), meaning
if \( f : X \to \mathbb{R} \) is cont. with respect to \( \sigma(X, X') \), then there is some unique
\( x' \in X' \) s.t. \( f(x) = \langle x, x' \rangle \) for all \( x \in X \).

A locally convex topology \( \tau \) on \( X \) is consistent (or compatible)
with the dual pair \( (X, X') \) if \( (X, \tau)' = X' \). 

Suprisingly, there are many stronger topologies which are consistent with \( (X, X') \) aside from \( \sigma(X, X') \).
Every topology consistent with a dual pair is Hausdorff.
All locally convex topologies consistent with a given dual pair have the same collection of closed convex sets.


Define the polar of a pairing

1. For a dual pair \((X, X')\), the absolute polar, or simply the polar, \(A^{\circ}\) of a nonempty subset \(A\) of \(X\), 
is the subset of \(X'\) defined by  
\[
A^{\circ} = \{ x' \in X' : \left| \langle x, x' \rangle \right| \leq 1 \text{ for all } x \in A \}.
\]
2. The bipolar of a subset \(A\) of \(X\) or \(X'\) is the set \((A^{\circ})^{\circ}\), written simply as \(A^{\circ \circ}\).  
3. The one-sided polar of \(A \subset X\), denoted \(A^{\ominus}\), is defined by  
\[
A^{\ominus} = \{ x' \in X' : \langle x, x' \rangle \leq 1 \text{ for all } x \in A \}.
\]

These definitions are analogous for \( X' \).

Let \(A, B\) be nonempty subsets of \(X\), and let \(\{A_i\}\) be a family of nonempty subsets of \(X\). Then:
1. If \(A \subset B\), then \(A^{\circ} \supset B^{\circ}\) and \(A^{\ominus} \supset B^{\ominus}\).
2. If \(\varepsilon \neq 0\), then \((\varepsilon A)^{\circ} = \frac{1}{\varepsilon} A^{\circ}\) and \((\varepsilon A)^{\ominus} = \frac{1}{\varepsilon} A^{\ominus}\).
3.
\[
\bigcap A_i^{\circ} = \left(\bigcup A_i\right)^{\circ}
\]
and
\[
\bigcap A_i^{\ominus} = \left(\bigcup A_i\right)^{\ominus}.
\]
4. The absolute polar \(A^{\circ}\) is nonempty, convex, circled, \(\sigma(X', X)\)-closed, and contains zero.
5. The one-sided polar \(A^{\ominus}\) is nonempty, convex, \(\sigma(X', X)\)-closed, and contains zero.
6. If \(A\) is absorbing, then both \(A^{\circ}\) and \(A^{\ominus}\) are \(\sigma(X', X)\)-bounded.
7. The set \(A\) is \(\sigma(X, X')\)-bounded if and only if \(A^{\circ}\) is absorbing.
The corresponding dual statements are true for subsets of \(X'\).


State the bipolar theorem and the polar version of Alaoglu’s compactness theorem

Bipolar Theorem  
Let \((X, X')\) be a dual pair, and let \(A\) be a nonempty subset of \(X\).

1. The bipolar \(A^{\circ\circ}\) is the convex circled \(\sigma(X, X')\)-closed hull of \(A\). Hence if \(A\) is convex, circled, and \(\sigma(X, X')\)-closed, then \(A = A^{\circ\circ}\).
2. The one-sided bipolar \(A^{\ominus\ominus}\) is the convex \(\sigma(X, X')\)-closed hull of \(A \cup \{0\}\). Hence if \(A\) is convex, \(\sigma(X, X')\)-closed, and contains zero, then \(A = A^{\ominus\ominus}\).
Corresponding results hold for subsets of \(X'\).


Let \( V \) be a neighborhood of zero for some locally convex topology on \( X \) consistent with the dual pair 
\( (X, X') \). Then its polar \( V^\circ \) is a \( w^\ast \)-compact subset of \( X' \).
Similarly, if \( W \) is a neighborhood of zero for some consistent locally convex
topology on \( X' \), then its polar \( W^\circ \) is a weakly compact subset of \( X \).

Relate the polar to the annhilator of a subspace.

Let \( (X, X') \) be a dual pair and let \( M \) be a linear subspace of \( X \).
Then \( M^{\bot} = M^{\ominus} = M^\circ \). If \( M \) is weakly closed, then \( M^{\bot \bot} = M \).
An analogous result holds for linear subspaces of \( X' \).

Let \( (X, X') \) be a dual pair and let
\( Y' \) be a linear subspace of \( X' \). Then the following are equivalent.
1. \( Y' \) is total. That is, \( Y' \) separates points of \( X \).
2. \( (Y')^{\bot} = \{0\} \).
3. \( Y' \) is weak* dense in \( X \).
The corresponding symmetric result is true for subspaces of \( X \).



Define Polar Topologies (\( \mathfrak{G} \)-topologies)

We start with an arbitrary nonempty \( \sigma(X', X) \)-bounded
\( A \subseteq X' \). Then 
\[
    q_A(x) = \sup \{ \|\langle x, x' \rangle : x' \in A \|\}  
\]
is a seminorm on \( X \). Further \( A^\circ = \{x \in X : q_A(x) \leq 1\} \)
and this is related to the Minkowski funtional via
\[
    q_A(x) = \inf \{ \alpha > 0 : x \in \alpha A^\circ \} = p_{A^\circ}(x)
\]
By the bipolar theorem we can also deduce \( q_A = q_{A^{\circ \circ}} \).

Now let \( S \) be a family of nonempty \( \sigma(X', X) \)-bounded subsets of \( X' \). 
The corresponding polar topology \( \mathfrak{S} \)-topology on \( X \) is the locally convex topology 
generated by the family of seminorms \( q_A : A \in S \). 
Equivalently, it is the topology generated by the neighborhood subbase (so all intersections of the following)
\( \{\varepsilon A^\circ : A \in S \text{ and } \varepsilon > 0\} \) 
at the origin.

The \( \mathfrak{S} \)-topology is Hausdorff if and only if the
span of the set \( \bigcup_{A \in \mathfrak{S}} A \) is \( \sigma(X', X) \)-dense in \( X' \). 
TFAE:
1. \( x_\alpha \to^{\mathfrak{S}} x \) in \( X \) 
2. \( q_A(x_\alpha − x) \to 0 \) for all \( A \in \mathfrak{S} \)
3. \( \{x_\alpha\} \) converges uniformly to \( x \) on each \( A \in \mathfrak{S} \).
this motivates to call the polar topology also the topology of uniform convergence on members of \( \mathfrak{S} \).

Remarkably, every consistent locally convex topology on X (or on \( X' \)) is an \( \mathfrak{S} \)-topology,
which is referred to as the Mackey–Arens Theorem.

The weak topology \( \sigma(X, X') \) is the coarsest locally convex topology on \( X \) consistent with 
\( (X, X') \). The finest consistent locally convex topology on \( X \) is by the \( \mathfrak{S} \)-topology for 
\( \mathfrak{S} \)consisting of all convex, circled, \( \sigma(X', X) \)-compact subsets of \( X' \). 
This is called the Mackey topology and denoted \( \tau(X, X') \).

Lemma (Mackey neighborhoods)  
If a nonempty subset \(K\) of \(X'\) is \(\sigma(X', X)\)-compact, 
then the one-sided polar \(K^{\ominus}\) is a convex \(\tau(X, X')\) (Mackey) neighborhood of zero in \(X\). 
Conversely, the one-sided polar \(V^{\ominus}\) of an arbitrary \(\tau(X, X')\)-neighborhood \(V\) of zero is 
nonempty, convex, and \(\sigma(X', X)\)-compact.

Note that the strong topology generated by all \( \sigma(X', X) \)-bounded subsets in \( X' \) is not always consistent with \( (X, X') \),
however for a normed space \( X \), the Mackey topology, the strong topology, and the norm topology are the same.
Its topological bidual is the traditional (algebraic) bidual, and is therefore relevant in questions of reflexivity for non normable TVS.


State the main differences for the topologies in the finite and infinite dimensional case 

For a normed space \( X \) tfae:
1. The vector space \( X \) is finite dimensional.
2. The weak and norm topologies on \( X \) coincide.
3. The weak topology on \( X \) is metrizable.
4. The weak topology is first countable.

The weak interior of every closed or open ball in an infinite dimensional normed space is empty.

In any infinite dimensional normed space, the closed unit sphere is weakly dense in the closed unit ball.
A normed space is separable if and only if the closed unit ball of its dual space is \( w^\ast \) -metrizable.

Banach-Alaoglu Theorem
If \( X \) is a normed space then the closed unit ball in the continuous dual space \( X' \)
(endowed with its usual operator norm) is compact with respect to the \( w^\ast \)-topology.

The dual \( X' \) of a Banach space \( X \) is separable if and only if the unit ball of \( X \) is weakly metrizable.
If the dual \( X' \) of a normed space \( X \) includes a countable total set, then every weakly compact subset of \( X \) is weakly metrizable.


State the standard collection of theorems for weak compactness in Banach spaces

Grothendiecks Theorem:
A subset \( A \) of a Banach space \( X \) is relatively weakly compact if and only if for each \( \varepsilon > 0 \) 
there exists a weakly compact set \( W \) such that \( A \subseteq W + \varepsilon U \), where \( U \) denotes the closed unit ball of \( X \).

Eberlein–Šmulian Theorem:
In the weak topology on a normed space, compactness and sequential compactness coincide. 
That is, a subset \(A\) of a normed space \(X\) is relatively weakly compact (respectively, weakly compact) 
if and only if every sequence in \(A\) has a weakly convergent subsequence in \(X\) (respectively, in \(A\)).

Krein–Šmulian Theorem:
In a Banach space, both the convex circled hull and the convex hull of a relatively weakly compact set are relatively weakly compact sets.

James' Theorem:
A nonempty weakly closed bounded subset of a Banach space is weakly compact if and only if 
every continuous linear functional attains a maximum on the set.


Bishop-Phelps
Let \(B \subseteq X\) be a bounded, closed, convex subset of a real Banach space \(X.\) 
Then the set of all cont. lin. \(f\) that achieve their supremum on \(B\) 
(meaning that there exists some \(b_0 \in B\) such that \(|f(b_0)| = \sup_{b \in B} |f(b)|\)) 
\[
\left\{f \in X^* : f \text{ attains its supremum on } B\right\} 
\]
is (dual norm)-dense in the \( X' \).
